\section{Maxima and truncations}

Truncating a potential requires stability under maxima. For a smooth background, this follows from the local psh maximum theorem by freezing the coefficients and then letting the curvature error tend to zero.

\begin{lemma}\label{lem:max-psh}
If $u,v\in\PSH(X,\theta)$, then $\max(u,v)\in\PSH(X,\theta)$. Consequently, if $\theta\geq0$ and $u\in\PSH(X,\theta)$, then $\max(u,-t)$ is $\theta$-plurisubharmonic for every real $t$.
\end{lemma}
\begin{proof}
The maximum is quasi-plurisubharmonic by a common local smooth correction. To verify its precise lower Hessian bound, fix a positive constant $(1,1)$-form $\kappa$ in a coordinate ball. For each $\varepsilon>0$, freeze $\theta$ on smaller neighborhoods to a constant form $\theta'$ with $-\varepsilon\kappa\leq\theta-\theta'\leq\varepsilon\kappa$. Write $\theta'+\varepsilon\kappa=\ddc r$ locally. Then $u+r$ and $v+r$ are plurisubharmonic, hence so is their maximum. It follows that
\[
 \theta+2\varepsilon\kappa+\ddc\max(u,v)\geq0
\]
on each such neighborhood. Let $\varepsilon\to0$. For the last assertion, constants belong to $\PSH(X,\theta)$ when $\theta\geq0$.
\end{proof}


\begin{lemma}\label{lem:strict-set-positive-volume}
If $a,b$ are quasi-plurisubharmonic functions and $\{a<b\}$ is nonempty at a point where both are finite, then this set has positive smooth volume.
\end{lemma}
\begin{proof}
Choose a point $x$ with $a(x)<b(x)$ and a number $c$ between them. Upper semicontinuity gives $a<c$ on a neighborhood $V$ of $x$. If $\{b>c\}\cap V$ had zero volume, then $b\leq c$ almost everywhere on $V$. The canonical quasi-plurisubharmonic representative, or the submean inequality after a smooth local correction, would imply $b(x)\leq c$, a contradiction. Thus $\{a<b\}$ has positive volume.
\end{proof}
